% Chapter 03 -- Content from chaptersnew, diagrams styled as chapters/.
\section{Chương 3: Phát triển và công cụ}
\ChapterDivider{Chương 3}{Phát triển và công cụ}{Từ mã nguồn đến quan sát được sự kiện thật}

\begin{frame}{Quy trình phát triển chương trình eBPF}
  \noindent{\small Ví dụ với C, Clang/LLVM và libbpf.}\par
  \begin{center}
    \begin{tikzpicture}
      \node[ebpfuser, text width=3.7cm, minimum height=1.9cm, align=left] (write) at (0,0)
        {\begin{minipage}[t][1.5cm][t]{3.7cm}\raggedright
         \textbf{1. Write}\\[0.13cm]
         Viết chương trình eBPF bằng C.
         \end{minipage}};
      \node[ebpfuser, text width=3.7cm, minimum height=1.9cm, align=left] (compile) at (4.7,0)
        {\begin{minipage}[t][1.5cm][t]{3.7cm}\raggedright
         \textbf{2. Compile}\\[0.13cm]
         Clang/LLVM tạo object chứa bytecode.
         \end{minipage}};
      \node[ebpfkernel, text width=3.7cm, minimum height=1.9cm, align=left] (load) at (9.4,0)
        {\begin{minipage}[t][1.5cm][t]{3.7cm}\raggedright
         \textbf{3. Load \& Verify}\\[0.13cm]
         Loader yêu cầu nạp, Verifier kiểm tra.
         \end{minipage}};
      \node[ebpfkernel, text width=3.7cm, minimum height=1.9cm, align=left] (attach) at (9.4,-2.5)
        {\begin{minipage}[t][1.5cm][t]{3.7cm}\raggedright
         \textbf{4. Attach}\\[0.13cm]
         Gắn program vào hook tương thích.
         \end{minipage}};
      \node[ebpfprogram, text width=3.7cm, minimum height=1.9cm, align=left] (execute) at (4.7,-2.5)
        {\begin{minipage}[t][1.5cm][t]{3.7cm}\raggedright
         \textbf{5. Execute}\\[0.13cm]
         Sự kiện phù hợp kích hoạt program.
         \end{minipage}};
      \node[ebpfmap, text width=3.7cm, minimum height=1.9cm, align=left] (collect) at (0,-2.5)
        {\begin{minipage}[t][1.5cm][t]{3.7cm}\raggedright
         \textbf{6. Collect}\\[0.13cm]
         User Space đọc map hoặc nhận sự kiện.
         \end{minipage}};
      \draw[arr] (write) -- (compile);
      \draw[arr] (compile) -- (load);
      \draw[arr] (load) -- (attach);
      \draw[arr] (attach) -- (execute);
      \draw[arr] (execute) -- (collect);
    \end{tikzpicture}
  \end{center}
  {\small \textbf{Load và Attach là hai bước riêng.} JIT là tùy chọn. Công cụ cấp cao có thể tự xử lý nhiều bước.}
  \note{Đây là quy trình ví dụ với C/libbpf. Loader yêu cầu kernel nạp chương trình; Verifier kiểm tra trong bước nạp. JIT có thể dịch bytecode sang mã máy khi được hỗ trợ và bật. bpftrace có thể quản lý nhiều bước thay người dùng. Nguồn: \url{https://docs.kernel.org/bpf/libbpf/libbpf_overview.html}; \url{https://docs.kernel.org/admin-guide/sysctl/net.html}.}
\end{frame}

\begin{frame}{Công cụ phát triển và sử dụng eBPF}
  \begin{center}
    \begin{tikzpicture}
      \node[ebpfuser, text width=6.0cm, minimum height=1.95cm, align=left] at (0,0)
        {\begin{minipage}[t][1.55cm][t]{6.0cm}\raggedright
         \textbf{bpftrace}\\[0.13cm]
         Tracing nhanh bằng cú pháp ngắn, phù hợp minh họa sự kiện.
         \end{minipage}};
      \node[ebpfuser, text width=6.0cm, minimum height=1.95cm, align=left] at (7.3,0)
        {\begin{minipage}[t][1.55cm][t]{6.0cm}\raggedright
         \textbf{bpftool}\\[0.13cm]
         Kiểm tra tính năng BPF của kernel. Quản lý program, map và link.
         \end{minipage}};
      \node[ebpfuser, text width=6.0cm, minimum height=1.95cm, align=left] at (0,-2.5)
        {\begin{minipage}[t][1.55cm][t]{6.0cm}\raggedright
         \textbf{libbpf}\\[0.13cm]
         Thư viện C hỗ trợ nạp, gắn và quản lý chương trình eBPF.
         \end{minipage}};
      \node[ebpfuser, text width=6.0cm, minimum height=1.95cm, align=left] at (7.3,-2.5)
        {\begin{minipage}[t][1.55cm][t]{6.0cm}\raggedright
         \textbf{BCC}\\[0.13cm]
         Thư viện và công cụ tracing để phân tích hành vi, hiệu năng hệ thống.
         \end{minipage}};
    \end{tikzpicture}
  \end{center}
  {\small \textbf{Thực hành:} Dùng \texttt{bpftool} kiểm tra môi trường và \texttt{bpftrace} theo dõi system calls, đếm sự kiện. Chưa cần tự viết loader bằng C.}
  \note{Vai trò và lựa chọn cho phần thực hành theo nội dung chaptersnew. Nguồn: \url{https://bpftrace.org/about}; \url{https://github.com/libbpf/bpftool/blob/main/docs/bpftool.rst}; \url{https://docs.kernel.org/bpf/libbpf/libbpf_overview.html}; \url{https://github.com/iovisor/bcc}.}
\end{frame}

\begin{frame}{Ví dụ: Theo dõi lời gọi \texttt{openat()}}
  \small
  Đếm lời gọi \texttt{openat()} yêu cầu mở tệp hoặc thư mục, thống kê theo \texttt{comm}.
  \begin{center}
    \begin{tikzpicture}
      \node[ebpfuser, text width=2.05cm, minimum height=2.4cm] (app) at (0,0)
        {\textbf{Application}\\[0.15cm]Gọi\\\texttt{openat()}};
      \node[ebpfkernel, text width=2.7cm, minimum height=2.4cm] (hook) at (3.0,0)
        {\textbf{Tracepoint}\\[0.15cm]\footnotesize\texttt{syscalls:}\\\texttt{sys\_enter\_openat}};
      \node[ebpfprogram, text width=2.15cm, minimum height=2.4cm] (prog) at (6.1,0)
        {\textbf{eBPF Program}\\[0.15cm]Lấy \texttt{comm}, cập nhật bộ đếm};
      \node[ebpfmap, text width=2.0cm, minimum height=2.4cm] (map) at (9.0,0)
        {\textbf{eBPF Map}\\[0.15cm]Số lần gọi theo\\\texttt{comm}};
      \node[ebpfuser, text width=2.15cm, minimum height=2.4cm] (user) at (11.9,0)
        {\textbf{bpftrace}\\[0.15cm]Đọc map, hiển thị thống kê};
      \draw[arr] (app) -- (hook);
      \draw[arr] (hook) -- (prog);
      \draw[arr] (prog) -- (map);
      \draw[arr] (map) -- (user);
    \end{tikzpicture}
  \end{center}
  \cue{1}{\texttt{comm} là tên tác vụ/tiến trình, không phải PID. Nhiều tiến trình có thể trùng tên.}
  \cue{2}{Sự kiện tại \textsl{entry} ghi nhận yêu cầu, chưa xác nhận mở tệp thành công.}
  \note{Sơ đồ minh họa cách bpftrace thu thập và hiển thị bộ đếm theo comm, không phải mã hay kết quả demo đã chạy. Tên tác vụ do kernel cung cấp có thể thay đổi và bị giới hạn độ dài. Tracepoint sys_enter_openat xảy ra khi bắt đầu lời gọi. Nguồn: \url{https://bpftrace.org/docs/release_024/language}; \url{https://man7.org/linux/man-pages/man5/proc_pid_comm.5.html}; \url{https://man7.org/linux/man-pages/man2/openat.2.html}.}
\end{frame}
